\subsection{Discrete Fourier Transforms and Circulant Matrices}

Throughout this subsection, $n\in\mathbb{Z}_{>0}$.

\begin{definition} \label{def:discrete-fourier-matrix} \hlt{Discrete Fourier Matrix}\\
Set $\omega_n\defeq e^{2\pi i/n}$, $F_n\defeq[\omega_n^{jk}]_{j,k=0}^{n-1}\in\operatorname{Mat}(n,n,\mathbb{C})$, and $Q_n\defeq n^{-1/2}F_n$.\\
For $\mathbf{c}\in\mathbb{C}^n$, $F_n\mathbf{c}$ is its unnormalized DFT and $Q_n\mathbf{c}$ its normalized DFT.
\end{definition}

\begin{definition} \label{def:circulant-matrix-cyclic-shift} \hlt{Circulant Matrix}\\
Let $(\mathbf{e}_0,\ldots,\mathbf{e}_{n-1})$ be the $0$-indexed standard basis of $\mathbb{C}^n$.\\
The \hlt{cyclic shift} is the matrix $P_n\in\operatorname{Mat}(n,n,\mathbb{C})$ defined by $(\forall\,j\in\{0,\ldots,n-1\})\,P_n\mathbf{e}_j=\mathbf{e}_{(j-1)\bmod n}$.\\
A matrix $C\in\operatorname{Mat}(n,n,\mathbb{C})$ is \hlt{circulant} iff $(\exists\,(a_0,\ldots,a_{n-1})\in\mathbb{C}^n)\,C=\sum\limits_{r=0}^{n-1}a_rP_n^r$.\\
If $(a_0,\ldots,a_{n-1})\in\mathbb{C}^n$ satisfies this identity, set $p(z)\defeq\sum\limits_{r=0}^{n-1}a_rz^r\in\mathbb{C}[z]$; then $C=p(P_n)$.\\
Here $p(P_n)$ denotes the polynomial evaluation from Definition \ref{def:polynomial-evaluation-on-square-matrix}.
\end{definition}

\begin{theorem} \label{thm:fourier-basis-diagonalizes-cyclic-circulants} \hlt{Fourier Diagonalization of Cyclic Shifts and Circulants}\\
Let $F_n,Q_n$ be as in Definition \ref{def:discrete-fourier-matrix} and $P_n$ as in Definition \ref{def:circulant-matrix-cyclic-shift}. Set $\Omega_n\defeq\operatorname{diag}(1,\omega_n,\ldots,\omega_n^{n-1})$.\\
For $(a_0,\ldots,a_{n-1})\in\mathbb{C}^n$, set $p(z)\defeq\sum\limits_{r=0}^{n-1}a_rz^r$ and $C\defeq p(P_n)$. Then:
\begin{enumerate}[label=(\roman*)]
\setlength{\itemsep}{0pt}
\item $F_n^*F_n=nI_n$, $Q_n^*Q_n=I_n$, and $Q_n$ is unitary;
\item $P_nF_n=F_n\Omega_n$ and $Q_n^*P_nQ_n=\Omega_n$;
\item $Q_n^*CQ_n=p(\Omega_n)=\operatorname{diag}\bigl(p(1),p(\omega_n),\ldots,p(\omega_n^{n-1})\bigr)$.
\end{enumerate}
\begin{proof}
For $k\in\{0,\ldots,n-1\}$, set $\mathbf{f}_k\defeq(\omega_n^{jk})_{j=0}^{n-1}$. Then $F_n=[\mathbf{f}_0\ \cdots\ \mathbf{f}_{n-1}]$ by Definition \ref{def:discrete-fourier-matrix}.
\begin{enumerate}[label=(\roman*)]
\setlength{\itemsep}{0pt}
\item For $k,\ell\in\{0,\ldots,n-1\}$, $(F_n^*F_n)_{k\ell}=\sum\limits_{r=0}^{n-1}\omega_n^{r(\ell-k)}$. This sum is $n$ if $k=\ell$.\\
If $k\neq\ell$, then $W\defeq\omega_n^{\ell-k}$ satisfies $W\neq1$ by $0<\abs{\ell-k}<n$ and $W^n=1$, so $\sum\limits_{r=0}^{n-1}W^r=\frac{1-W^n}{1-W}=0$.\\
Thus $F_n^*F_n=nI_n$ and $Q_n^*Q_n=n^{-1}F_n^*F_n=I_n$. Hence $Q_n$ is unitary by Theorem \ref{thm:unitary-matrix-equivalent-conditions}(i).
\item For $k,j\in\{0,\ldots,n-1\}$, $(P_n\mathbf{f}_k)_j=(\mathbf{f}_k)_{(j+1)\bmod n}=\omega_n^k(\mathbf{f}_k)_j$ by $\omega_n^{nk}=1$.\\
Hence $P_nF_n=F_n\Omega_n$ and $P_nQ_n=Q_n\Omega_n$ by column comparison. Therefore $Q_n^*P_nQ_n=\Omega_n$ by (i).
\item Set $D\defeq\operatorname{diag}\bigl(p(1),p(\omega_n),\ldots,p(\omega_n^{n-1})\bigr)$. By (i),(ii), $P_n=Q_n\Omega_nQ_n^*$.\\
The identity $C=Q_nDQ_n^*$ holds by Theorem \ref{thm:polynomial-functional-calculus-diagonalizable-matrix} with $(A,S,\Lambda)=(P_n,Q_n,\Omega_n)$.\\
The identity $p(\Omega_n)=D$ holds by the same theorem with $(A,S,\Lambda)=(\Omega_n,I_n,\Omega_n)$. Hence $Q_n^*CQ_n=D=p(\Omega_n)$ by (i).
\end{enumerate}
\hlt{Structural Commentary}: Root orthogonality proves unitarity; shift eigenvectors transfer polynomial evaluation to the diagonal.\\
\hlt{Key Insight}: The Fourier modes form an orthonormal eigenbasis of the cyclic shift.
\end{proof}
\end{theorem}

Figure \ref{fig:circulant-fourier-diagonalization} summarizes Theorem \ref{thm:fourier-basis-diagonalizes-cyclic-circulants}. Horizontal arrows conjugate by $Q_n$, and vertical arrows evaluate $p$.

\begin{illustrationfigure}[H]
  \begin{tikzpicture}[
  >=Stealth,
  x=0.92cm,
  y=0.9cm,
  every node/.style={font=\small},
  box/.style={draw=gray!55, rounded corners=3pt, fill=gray!6, align=center, minimum width=2.15cm, minimum height=0.8cm},
  result/.style={box, fill=green!7},
  arrow/.style={->, line width=0.85pt, blue!70!black},
  lbl/.style={fill=white, text=gray!70!black, inner xsep=2pt, inner ysep=1pt}
]
  \node[box] (shift) at (0,1.1) {$P_n$};
  \node[box] (diagshift) at (4.15,1.1) {$\Omega_n$};
  \node[box] (circ) at (0,-1.1) {$C=p(P_n)$};
  \node[result] (diagcirc) at (4.15,-1.1) {$p(\Omega_n)$};

  \draw[arrow] (shift) -- node[lbl, above] {$Q_n^*(\,\cdot\,)Q_n$} (diagshift);
  \draw[arrow] (circ) -- node[lbl, below] {$Q_n^*(\,\cdot\,)Q_n$} (diagcirc);
  \draw[arrow] (shift) -- node[lbl, left] {$p(\cdot)$} (circ);
  \draw[arrow] (diagshift) -- node[lbl, right] {$p(\cdot)$} (diagcirc);
\end{tikzpicture}

  \caption{Fourier conjugation commutes with polynomial evaluation: $P_n$ maps to $\Omega_n$ and $C=p(P_n)$ to $p(\Omega_n)$.}
  \label{fig:circulant-fourier-diagonalization}
\end{illustrationfigure}

\begin{proposition} \label{prop:fourier-matrix-inverse} \hlt{Inverse of the Fourier Matrix}\\
The Fourier matrix satisfies $F_n^{-1}=\frac{1}{n}F_n^*$.
\begin{proof}
The identity $\bigl(\frac{1}{n}F_n^*\bigr)F_n=I_n$ holds by Theorem \ref{thm:fourier-basis-diagonalizes-cyclic-circulants}(i). Hence $F_n^{-1}=\frac{1}{n}F_n^*$ by Theorem \ref{thm:square-one-sided-inverse-is-two-sided}.\\
\hlt{Structural Commentary}: Fourier orthogonality supplies a left inverse; squareness identifies the two-sided inverse.\\
\hlt{Key Insight}: Normalization turns conjugate transpose into exact inversion.
\end{proof}
\end{proposition}
